\documentclass[../main]{subfiles}

\begin{document}

\chapter{Segundo Principio de la Termodinámica: Ciclo de Carnot}

\section{Teoría}

(Continúa la sección de teoría del ciclo de Carnot...)

\section{Fórmulas}

(Continúa la sección de fórmulas del ciclo de Carnot...)

\section{Cuestionario}

(Continúa la sección de cuestionario del ciclo de Carnot...)

\section{Problemas}

1. Un motor térmico opera entre un reservorio caliente a 500 °C y un reservorio frío a 25 °C. Si el motor absorbe 1000 J de calor del reservorio caliente, ¿cuál es la cantidad de trabajo que realiza y la cantidad de calor que libera al reservorio frío?

**Solución:**

* Calor absorbido: $Q_H = 1000 J$
* Temperatura del reservorio caliente: $T_H = 500 °C = 773 K$
* Temperatura del reservorio frío: $T_C = 25 °C = 298 K$

* Trabajo realizado:

$$
W = \eta Q_H = \frac{T_H - T_C}{T_H} Q_H = \frac{773 K - 298 K}{773 K} \cdot 1000 J = 610.71 J
$$

* Calor liberado:

$$
Q_C = Q_H - W = 1000 J - 610.71 J = 389.29 J
$$

2. Una central eléctrica opera con un ciclo de Carnot cuya eficiencia es del 60%. Si la temperatura del reservorio frío es de 20 °C, ¿cuál es la temperatura máxima que puede alcanzar el reservorio caliente?

**Solución:**

* Eficiencia: $\eta = 60\% = 0.6$
* Temperatura del reservorio frío: $T_C = 20 °C = 293 K$

* Temperatura del reservorio caliente:

$$
T_H = \frac{T_C}{1 - \eta} = \frac{293 K}{1 - 0.6} = 732.5 K
$$

3. Un motor de combustión interna tiene una eficiencia del 30%. Si el motor quema 1 kg de gasolina, que libera 40 000 kJ de energía, ¿cuál es la potencia útil del motor?

**Solución:**

* Eficiencia: $\eta = 30\% = 0.3$
* Energía liberada por la combustión: $Q = 40 000 kJ = 40 000 000 J$

* Energía útil:

$$
W_u = \eta Q = 0.3 \cdot 40 000 000 J = 12 000 000 J
$$

* Potencia útil:

$$
P = \frac{W_u}{t}
$$

Donde $t$ es el tiempo en que se consume el combustible. Para este problema, no se ha proporcionado información sobre el tiempo de consumo, por lo que no se puede calcular la potencia útil.

**Bibliografía:**

* Termodinámica - Stevenazzi
* Curso de Termodinámica - Facorro Ruiz
* Física - Steinway
* Física - Sears Semansky

\end{document}
